% ============================================================
%  S&P 500 Stock Price Forecasting using Deep Learning
%  COMS 4030A / COMS 7047A -- Adaptive Computation and Machine Learning
%  University of Cape Town -- May 2026
%
%  Template: colacl (ACML / ACL conference paper format)
%  Requires colacl.sty in the same directory (download from the
%  shared Overleaf project).
% ============================================================
\documentclass[11pt]{article}
\usepackage{colacl}
\usepackage{graphicx}
\usepackage{booktabs}
\usepackage{tabularx}
\usepackage{array}
\usepackage{amsmath}
\usepackage{float}
\usepackage{enumitem}
\usepackage[colorlinks=true,linkcolor=black,urlcolor=black,citecolor=black]{hyperref}
\sloppy

% Column types for tables
\newcolumntype{L}[1]{>{\raggedright\arraybackslash}p{#1}}
\newcolumntype{C}[1]{>{\centering\arraybackslash}p{#1}}

\title{S\&P 500 Stock Price Forecasting using Deep Learning}

\author{
  Zeeshaan Ebrahim \quad {\small (2543004)} \And
  Tazeem Tayob \quad {\small (2659956)}
  \\[6pt]
  {\small COMS 4030A / COMS 7047A --- Adaptive Computation and Machine Learning} \\
  {\small University of Cape Town, May 2026}
}

\begin{document}
\maketitle

% ============================================================
\begin{abstract}
This report presents a comparative study of four deep learning
architectures---Stacked LSTM, Bidirectional LSTM with Bahdanau
Attention, CNN-LSTM Hybrid, and a custom Transformer Encoder---for
next-day adjusted closing price forecasting across nine S\&P~500
technology stocks (2010--2026). Models were trained on a unified
dataset with 28 engineered features per timestep. Controlled
experiments investigated the effects of window size, feature richness,
dropout, prediction horizon, and L2 regularisation. The CNN-LSTM
Hybrid achieved the best performance (RMSE \$41.83, MAE \$26.79,
MAPE~10.28\%), demonstrating that causal convolutional feature
extraction combined with recurrent long-range modelling is
particularly effective for financial time-series forecasting.
\end{abstract}

% ============================================================
\section{Dataset Description \& Objectives}

\textbf{Source:} Kaggle --- S\&P 500 Top 10 Stocks OHLCV (2010--2026).
The dataset contains split-adjusted daily trading data for nine S\&P~500
technology companies, spanning over 15 years of market history.

\begin{table}[h]
\caption{Dataset Columns}
\label{tab:dataset-cols}
\begin{tabularx}{\linewidth}{L{2.8cm} X}
\toprule
\textbf{Column} & \textbf{Description} \\
\midrule
Date        & Trading date (YYYY-MM-DD) \\
Ticker      & Stock symbol (9 unique tickers) \\
Open        & Opening price (split-adjusted) \\
High        & Daily high price \\
Low         & Daily low price \\
Close       & Closing price \\
Adj\_Close  & Adjusted closing price (primary target) \\
Volume      & Number of shares traded \\
\bottomrule
\end{tabularx}
\end{table}

\noindent\textbf{Tickers:} AAPL, AMZN, AVGO, GOOG, GOOGL, META,
MSFT, NVDA, TSLA

\noindent\textbf{Dataset statistics:} 35,765 rows $|$ 9 stocks $|$
Date range: 2010-01-04 to 2026-02-13 $|$ Missing values: 0

\subsection{Objectives}

\begin{itemize}
  \item Build a single unified deep learning model to forecast the
        next trading day's adjusted closing price for any of the nine
        stocks simultaneously.
  \item Compare four architectures: Stacked LSTM, Bidirectional LSTM
        with Attention, CNN-LSTM Hybrid, and custom Transformer Encoder.
  \item Investigate how window size, feature richness, dropout,
        prediction horizon, and L2 regularisation affect forecast
        accuracy.
\end{itemize}

% ============================================================
\section{Preprocessing \& Data Splits}

\subsection{Feature Engineering}

17 additional technical indicators were computed per stock using only
information available up to time $t$ (no lookahead bias). All rolling
calculations were performed per-ticker to prevent cross-stock
contamination.

\begin{table}[h]
\caption{Engineered Features}
\label{tab:features}
\begin{tabularx}{\linewidth}{L{3.2cm} X}
\toprule
\textbf{Feature Group} & \textbf{Features} \\
\midrule
Trend (SMA)   & SMA\_5, SMA\_20, SMA\_50 \\
Trend (EMA)   & EMA\_12, EMA\_26, MACD \\
Momentum      & RSI\_14 \\
Volatility    & BB\_upper, BB\_lower, BB\_bandwidth \\
Volume        & Volume\_MA\_10 \\
Returns       & Daily\_Return, Volatility\_10 \\
Identity      & Ticker one-hot encoding (9 dimensions) \\
\bottomrule
\end{tabularx}
\end{table}

\noindent\textbf{Total: 28 features per timestep.}

\subsection{Normalisation}

MinMaxScaler was fitted exclusively on the training slice of each stock
to prevent data leakage. The same fitted scaler was then applied to the
validation and test slices. Scalers were stored per-ticker to ensure
correct inverse transformation during evaluation.

\subsection{Data Splits}

\begin{table}[h]
\caption{Data Splits}
\label{tab:splits}
\begin{tabularx}{\linewidth}{L{1.8cm} L{3.4cm} C{2.0cm} X}
\toprule
\textbf{Split} & \textbf{Date Range} & \textbf{Seqs.} & \textbf{Purpose} \\
\midrule
Train      & 2010-01-04 to 2022-12-31 & 27,746 & Fit models and learn weights \\
Validation & 2023-01-01 to 2024-06-30 & 2,826  & Hyperparameter tuning and early stopping \\
Test       & 2024-07-01 to 2026-02-13 & 3,132  & Final evaluation (never seen during training) \\
\bottomrule
\end{tabularx}
\end{table}

\subsection{Sequence Creation}

A sliding window of 60 trading days ($\approx$3 months) was used.
For each sample, $X =$ features for days $[t-60,\,t-1]$ and $y =$
normalised Adj\_Close at day $t$. Sequences from all nine stocks were
stacked into a single shared dataset.

% ============================================================
\section{Model Selection \& Implementation}

All models were implemented in PyTorch 2.2.2 and trained with the
Adam optimiser ($\text{lr}=10^{-3}$), MSE loss, gradient clipping
(max norm $= 1.0$), EarlyStopping (patience $= 10$), and
ReduceLROnPlateau (factor $= 0.5$, patience $= 5$).

\subsection{Model 1 --- Stacked LSTM (Baseline)}

A two-layer stacked LSTM with dropout regularisation serves as the
baseline.

\noindent\textbf{Architecture:}
$\text{Input}(60,28) \to \text{LSTM}(128,\;\text{return\_seq}=\text{True})
\to \text{Dropout}(0.2) \to \text{LSTM}(64) \to \text{Dropout}(0.2)
\to \text{Linear}(32) \to \text{ReLU} \to \text{Linear}(1)$

\noindent\textbf{Parameters:} 132,673 \quad $|$ \quad
\textbf{Best validation loss:} 0.070183

\subsection{Model 2 --- Bidirectional LSTM with Bahdanau Attention}

Bidirectional LSTMs capture both forward and backward temporal context
within the 60-day window. A custom Bahdanau-style attention mechanism
computes a weighted sum of hidden states, allowing the model to focus
on the most predictive timesteps. The final context vector is passed
through a fully connected head.

\noindent\textbf{Parameters:} 330,946 \quad $|$ \quad
\textbf{Best validation loss:} 0.115797

\subsection{Model 3 --- CNN-LSTM Hybrid}

Temporal causal convolutions act as local pattern extractors detecting
short-term momentum shifts, while the LSTM captures long-range
sequential dependencies. Causal padding (left-padding only) prevents
any future information leakage into the convolutional feature maps.

\noindent\textbf{Parameters:} 51,169 \quad $|$ \quad
\textbf{Best validation loss:} 0.012034

\subsection{Model 4 --- Custom Transformer Encoder}

Built entirely from scratch using PyTorch primitives following the
Vaswani et al.\ (2017) encoder design with multi-head self-attention,
position-wise feed-forward networks, layer normalisation, and residual
connections. Sinusoidal positional encodings were applied to the input
embeddings.

\noindent\textbf{Parameters:} 76,865 \quad $|$ \quad
\textbf{Best validation loss:} 0.122278

% ============================================================
\section{Experiments \& Fine-Tuning}

Five controlled experiments were conducted using the LSTM as a
baseline (except Experiment 5 which used the Transformer).

\subsection{Experiment 1 --- Window Size}

The sequence length was varied across 20, 30, and 60 trading days to
assess the impact of lookback horizon on forecast accuracy.

\begin{table}[h]
\caption{Window Size vs.\ Validation RMSE}
\label{tab:window}
\begin{tabularx}{\linewidth}{C{4.0cm} C{4.0cm}}
\toprule
\textbf{Window Size} & \textbf{Val RMSE} \\
\midrule
20 days & 0.1599 \\
30 days & 0.1952 \\
\textbf{60 days (best)} & \textbf{0.0749} \\
\bottomrule
\end{tabularx}
\end{table}

\noindent\textbf{Result:} A 60-day window significantly outperformed
shorter windows, capturing quarterly price trends and momentum cycles
invisible at shorter horizons.

\subsection{Experiment 2 --- Feature Ablation}

OHLCV-only (6 features) was compared against the full technical
indicator set (28 features).

\begin{table}[h]
\caption{Feature Set vs.\ Validation RMSE}
\label{tab:feat-ablation}
\begin{tabularx}{\linewidth}{X C{3.5cm}}
\toprule
\textbf{Feature Set} & \textbf{Val RMSE} \\
\midrule
OHLCV-only (6 features) & Higher RMSE \\
\textbf{Full features with indicators (28)} & \textbf{Lower RMSE (winner)} \\
\bottomrule
\end{tabularx}
\end{table}

\noindent\textbf{Result:} Technical indicators (MACD, RSI, Bollinger
Bands) provided meaningful additional signal, reducing RMSE compared
to raw price and volume data alone.

\subsection{Experiment 3 --- Dropout Rate}

Three dropout rates were tested on the LSTM.

\begin{table}[h]
\caption{Dropout Rate vs.\ Validation RMSE}
\label{tab:dropout}
\begin{tabularx}{\linewidth}{C{4.0cm} C{4.0cm}}
\toprule
\textbf{Dropout Rate} & \textbf{Val RMSE} \\
\midrule
\textbf{0.1 (best)} & \textbf{0.1325} \\
0.2 & 0.2322 \\
0.3 & 0.3171 \\
\bottomrule
\end{tabularx}
\end{table}

\noindent\textbf{Result:} A dropout rate of 0.1 achieved the lowest
validation RMSE. Higher dropout rates caused underfitting.

\subsection{Experiment 4 --- Prediction Horizon}

1-day and 5-day ahead forecasting were compared on the LSTM.

\begin{table}[h]
\caption{Prediction Horizon vs.\ Validation RMSE}
\label{tab:horizon}
\begin{tabularx}{\linewidth}{C{4.0cm} C{4.0cm}}
\toprule
\textbf{Horizon} & \textbf{Val RMSE} \\
\midrule
\textbf{1 day ahead (winner)} & \textbf{Lower RMSE} \\
5 days ahead & Higher RMSE ($\approx$doubled) \\
\bottomrule
\end{tabularx}
\end{table}

\noindent\textbf{Result:} The 5-day horizon was significantly harder,
with RMSE roughly doubling, indicating rapid error accumulation over
multiple timesteps.

\subsection{Experiment 5 --- L2 Regularisation (Transformer)}

Weight decay was applied to the Transformer optimiser across three
settings.

\begin{table}[h]
\caption{Weight Decay vs.\ Validation RMSE (Transformer)}
\label{tab:l2}
\begin{tabularx}{\linewidth}{C{4.0cm} C{4.0cm}}
\toprule
\textbf{Weight Decay} & \textbf{Val RMSE} \\
\midrule
0 (none) & 0.3389 \\
\textbf{$10^{-4}$ (best)} & \textbf{0.3250} \\
$10^{-3}$ & 0.3484 \\
\bottomrule
\end{tabularx}
\end{table}

\noindent\textbf{Result:} Mild L2 regularisation ($10^{-4}$) slightly
improved generalisation. Strong regularisation ($10^{-3}$) caused
underfitting.

% ============================================================
\section{Training Visualisations}

\begin{figure}[h]
  \centering
  \includegraphics[width=\linewidth]{fig1}
  \caption{Normalised adjusted closing prices for all 9 stocks
    (2010--2026).}
  \label{fig:norm-prices}
\end{figure}

\begin{figure}[h]
  \centering
  \includegraphics[width=\linewidth]{fig2}
  \caption{Cross-stock return correlation heatmap.}
  \label{fig:corr}
\end{figure}

\begin{figure}[h]
  \centering
  \includegraphics[width=\linewidth]{fig3}
  \caption{Daily trading volume distribution by stock.}
  \label{fig:volume}
\end{figure}

\begin{figure}[h]
  \centering
  \includegraphics[width=\linewidth]{fig4}
  \caption{Training and validation MSE loss curves for all four models.}
  \label{fig:loss}
\end{figure}

\begin{figure}[h]
  \centering
  \includegraphics[width=\linewidth]{fig5}
  \caption{Experiment results: window size, feature ablation, and
    dropout rate.}
  \label{fig:exp135}
\end{figure}

\begin{figure}[h]
  \centering
  \includegraphics[width=\linewidth]{fig6}
  \caption{Experiment results: prediction horizon and L2
    regularisation.}
  \label{fig:exp45}
\end{figure}

\begin{figure}[h]
  \centering
  \includegraphics[width=\linewidth]{fig7}
  \caption{BiLSTM attention weight heatmap over the 60-day lookback
    window.}
  \label{fig:attention}
\end{figure}

% ============================================================
\section{Test Results \& Analysis}

All predictions were inverse-transformed from normalised $[0,1]$ space
back to actual dollar prices before computing metrics.

\begin{table}[h]
\caption{Test Set Performance Across All Models}
\label{tab:results}
\begin{tabularx}{\linewidth}{L{2.8cm} C{1.6cm} C{1.6cm} C{1.8cm} C{2.0cm}}
\toprule
\textbf{Model} &
\textbf{RMSE (\$)} &
\textbf{MAE (\$)} &
\textbf{MAPE (\%)} &
\textbf{Dir.\ Acc.\ (\%)} \\
\midrule
LSTM                 & 79.953  & 55.399 & 19.857 & 49.38 \\
BiLSTM + Attn.       & 98.377  & 69.123 & 23.934 & 50.40 \\
\textbf{CNN-LSTM}    & \textbf{41.827} & \textbf{26.785} & \textbf{10.278} & \textbf{51.17} \\
Transformer          & 101.240 & 71.212 & 24.502 & 50.46 \\
\bottomrule
\end{tabularx}
\end{table}

\noindent\textbf{Best model:} CNN-LSTM with RMSE \$41.83, MAE \$26.79,
MAPE~10.28\%, Directional Accuracy 51.17\%.

\begin{figure}[h]
  \centering
  \includegraphics[width=\linewidth]{fig8}
  \caption{Test set performance comparison --- RMSE, MAE, MAPE across
    all models.}
  \label{fig:perf}
\end{figure}

\begin{figure}[h]
  \centering
  \includegraphics[width=\linewidth]{fig9}
  \caption{Actual vs.\ predicted Adj\_Close for AAPL, NVDA, TSLA
    (test set).}
  \label{fig:actual-pred}
\end{figure}

\begin{figure}[h]
  \centering
  \includegraphics[width=\linewidth]{fig10}
  \caption{Residual error distribution per stock --- CNN-LSTM (best
    model).}
  \label{fig:residuals}
\end{figure}

\subsection{Analysis}

The CNN-LSTM hybrid achieved the best performance across all metrics
by a significant margin. Its RMSE of \$41.83 was roughly half that of
the LSTM (\$79.95) and less than half that of the Transformer
(\$101.24). The local pattern extraction performed by the causal
convolutional layers proved highly effective at capturing short-term
momentum shifts before they are passed to the LSTM for longer-range
modelling.

The Bidirectional LSTM with Attention underperformed relative to
expectations. While attention mechanisms are theoretically beneficial,
the added model complexity (330K parameters, $2.5\times$ more than the
LSTM baseline) led to overfitting despite dropout regularisation.

The Transformer Encoder was the weakest performer (\$101.24 RMSE).
Without pre-training on large datasets, Transformers tend to require
substantially more data to generalise---the approximately 27,000
training sequences available were insufficient to exploit the
architecture's full capacity.

All four models achieved directional accuracy close to 50\%, slightly
better than random chance. This is consistent with the Efficient Market
Hypothesis, which posits that short-term price direction is inherently
difficult to predict from historical price data alone.

Residual distributions for all stocks are approximately centred around
zero, indicating no systematic bias. However, residuals are wider for
high-volatility stocks (NVDA, TSLA) compared to more stable equities
(AAPL, MSFT), reflecting the greater difficulty of forecasting highly
volatile assets.

% ============================================================
\section{Conclusions}

Four deep learning architectures were trained on a unified dataset of
nine S\&P~500 technology stocks to forecast next-day adjusted closing
prices. The CNN-LSTM hybrid was the clear best performer, achieving the
lowest RMSE, MAE, and MAPE across all nine stocks.

\subsection{What Worked}

\begin{itemize}
  \item Engineering technical indicators (MACD, RSI, Bollinger Bands)
        reduced RMSE compared to raw OHLCV data alone.
  \item A 60-day lookback window significantly outperformed 20- and
        30-day windows, capturing quarterly trend cycles.
  \item Causal convolutional feature extraction (CNN-LSTM) was more
        effective than pure recurrent or attention-based approaches.
  \item ReduceLROnPlateau scheduling and early stopping prevented
        overfitting and reduced training time.
\end{itemize}

\subsection{What Did Not Work}

\begin{itemize}
  \item The 5-day prediction horizon was significantly harder, with
        RMSE roughly doubling.
  \item High dropout rates (0.3) degraded performance, suggesting
        underfitting beyond this threshold.
  \item Strong L2 regularisation (weight\_decay $= 10^{-3}$) caused
        underfitting in the Transformer.
  \item The Bidirectional LSTM with Attention overfitted despite
        dropout regularisation.
\end{itemize}

\subsection{Limitations}

\begin{itemize}
  \item No fundamental data (earnings, macroeconomic indicators) were
        used; models rely solely on technical price history.
  \item Financial time-series are non-stationary; concept drift events
        (e.g.\ market crashes) pose challenges not addressed here.
  \item Directional accuracy near 50\% limits practical trading
        utility.
\end{itemize}

\subsection{Future Work}

\begin{itemize}
  \item \textbf{Cross-stock attention:} allow each stock to attend to
        signals from the other eight.
  \item \textbf{Multi-step forecasting:} recursive forecasting with
        uncertainty quantification via Monte Carlo Dropout.
  \item \textbf{Ensemble methods:} combine all four architectures for
        lower-variance predictions.
  \item \textbf{Richer features:} incorporate fundamental and
        macroeconomic indicators.
\end{itemize}

\end{document}
